\begin{center}
\begin{tikzpicture}[scale=1.25, line cap=round, line join=round]
  \definecolor{trigblue}{RGB}{30,110,190}
  \coordinate (A) at (0,0);
  \coordinate (B) at (3.4,0);
  \coordinate (C) at (3.4,2.0);
  \draw[very thick] (A)--(B)--(C)--cycle;
  \draw[line width=2.2pt,trigblue] (A)--(B);
  \draw[line width=2.2pt,gray!65] (A)--(C);
  \draw (3.1,0)--(3.1,0.3)--(3.4,0.3);
  \draw[trigblue] (0.55,0) arc[start angle=0,end angle=30.5,radius=0.55];
  \node[trigblue] at (0.78,0.22) {$\theta$};
  \node[below,trigblue] at ($(A)!0.5!(B)$) {adjacent};
  \node[above left] at ($(A)!0.5!(C)$) {hypotenuse};
  \node[below] at (1.7,-0.55) {$\cos\theta=\text{adjacent}/\text{hypotenuse}$};
\end{tikzpicture}
\end{center}


\begin{definitionbox}{Definition}
\begin{definition}[Right-triangle cosine]
\label{def:right-triangle-cosine}
Let $T$ be a right triangle and let $\theta$ be one of its acute angles. Define
\[
\cos\theta=\frac{\operatorname{adj}(T,\theta)}{\operatorname{hyp}(T)}.
\]
\end{definition}
\end{definitionbox}

\begin{center}
\begin{tikzpicture}[scale=1.25, line cap=round, line join=round]
  \definecolor{trigblue}{RGB}{30,110,190}
  \coordinate (A) at (0,0);
  \coordinate (B) at (3.4,0);
  \coordinate (C) at (3.4,2.0);
  \draw[very thick] (A)--(B)--(C)--cycle;
  \draw[line width=2.1pt,trigblue] (A)--(B);
  \draw[line width=2.1pt,gray!65] (A)--(C);
  \draw (3.1,0)--(3.1,0.3)--(3.4,0.3);
  \draw[trigblue] (0.55,0) arc[start angle=0,end angle=30.5,radius=0.55];
  \node[trigblue] at (0.78,0.22) {$\theta$};
  \node[below,trigblue] at ($(A)!0.5!(B)$) {$\operatorname{adj}(T,\theta)$};
  \node[above left] at ($(A)!0.5!(C)$) {$\operatorname{hyp}(T)$};
  \node[below] at (1.7,-0.55) {$\cos\theta=\operatorname{adj}(T,\theta)/\operatorname{hyp}(T)$};
\end{tikzpicture}
\end{center}


\begin{definitionbox}{Definition}
\begin{definition}[Acute-angle cosine]
\label{def:acute-angle-cosine}
For an acute Euclidean angle $\theta$, choose any right triangle $T$ having
$\theta$ as an acute angle and set
\[
  \cos\theta=\frac{\operatorname{adj}(T,\theta)}{\operatorname{hyp}(T)}.
\]
\end{definition}
\end{definitionbox}

\begin{center}
\begin{tikzpicture}[scale=1.55, line cap=round, line join=round]
  \definecolor{trigblue}{RGB}{30,110,190}
  \coordinate (O) at (0,0);
  \coordinate (P) at (35:1);
  \coordinate (Q) at ({cos(35)},0);
  \draw[gray!55] (O) circle (1);
  \draw[->] (-1.25,0)--(1.35,0) node[right] {$x$};
  \draw[->] (0,-1.15)--(0,1.25) node[above] {$y$};
  \draw[very thick, trigblue] (O)--(P);
  \draw[very thick] (Q)--(P);
  \draw[very thick] (O)--(Q);
  \draw (Q)+(-0.12,0)--++(0,0.12)--++(0.12,0);
  \draw[trigblue] (0.32,0) arc[start angle=0,end angle=35,radius=0.32];
  \node[trigblue] at (0.32,0.10) {$\theta$};
  \fill (O) circle (0.018) node[below left] {$O$};
  \fill (P) circle (0.022) node[above right] {$P(\theta)=(\cos\theta,\sin\theta)$};
  \node[below] at ($(O)!0.5!(Q)$) {$\cos\theta$};
  \node[right] at ($(Q)!0.5!(P)$) {$\sin\theta$};
\end{tikzpicture}
\end{center}


\begin{definitionbox}{Definition}
\begin{definition}[Unit-circle cosine]
\label{def:unit-circle-cosine}
Let $\theta$ be an angle in standard position, and let $P(\theta)=(x_\theta,
y_\theta)$ be the point where its terminal ray meets the unit circle. Define
\[
  \cos\theta=x_\theta.
\]
\end{definition}
\end{definitionbox}
